\begin{enumerate}
    \item Les forces appliquées à la bille du pendule sont~:
          \begin{itemize}
              \item son poids $\vec{P}$ (vertical, vers le bas, appliqué en son
                    centre d'inertie $G$)
              \item la tension $\vec{T}$ du fil (dirigée suivant le fil, vers le
                    point $O$ et appliquée au point d'attache).
          \end{itemize}

          \begin{figure}[H]
              \centering
              \begin{tikzpicture}
                  \def\th{30}
                  \def\tho{60}
                  \newdimen\el \el=4cm
                  \node[draw,thick,pattern=north east lines,
                        minimum width=1cm,minimum height=5mm] (S) {};
                  \coordinate (O) at (S.south);
                  \path[draw,fill=DarkTurquoise]
                      (O) -- ([yshift=-1cm]O) arc (270:270+\tho:1cm) --
                      cycle;
                  \node[circle,fill,inner sep=1.1pt,
                        label={[anchor=north east,inner sep=1.8pt]:$O$}]
                        at (O) {};
                  \draw[densely dashed]
                      (O) +(270-\tho:\el) arc (270-\tho:270+\tho:\el);
                  \foreach \p/\a in {A/\tho,G/\th,B/0}{%
                      \node[circle,ball color=black!60,inner sep=0pt,
                            minimum size=4mm,
                            label={[inner sep=1.6pt]-90+\a:$\p$}]
                            (\p) at ($(O)+(270+\a:\el)$) {};
                      \fill (\p.center) circle (1pt);
                      \draw[thick] (O) -- (\p);
                  }
                  \node[anchor=180+\tho,inner sep=1.4pt]
                      at ($(O)!.5!(A.center)$) {$\ell$};
                %node[pos=0.9,below] {$\theta_{0}$};
                  \draw[densely dashed] (A) -- (A -| O) coordinate (H)
                      node[anchor=south west,inner sep=1.2pt] {$H$};
                  \draw[thick,<->] (H) +(-3mm,0) coordinate(tmp) --
                      node[left] {$h$} (tmp |- B.center);
                  \draw[thick] (H) +(0,-6pt) -| +(6pt,0);
                  \draw[->,very thick,red] (G) -- ++(90+\th:0.5\el)
                      node[pos=0.6,right=1mm,inner sep=1.0pt] {$\vec{T}$};
                  \draw[->,very thick,blue] (G.center) -- ++(270:0.3\el)
                      node[pos=0.8,left] {$\vec{P}$};
              \end{tikzpicture}
          \end{figure}
          
    \item 
          \begin{itemize}
              \item Travail de la tension $\vec{T}$~:
                    
                    La trajectoire du point d'application de $\vec{T}$ est un
                    cercle de centre $O$. $\vec{T}$ est suivant le rayon de ce
                    cercle et un déplacement élémentaire $\vv{MM}^{\prime}$
                    coïncide avec la tangente au cercle~; la force $\vec{T}$ est
                    donc constamment perpendiculaire au déplacement et
                    $W(\vec{T})=0$.
              \item Travail du poids $\vec{P}$~:

                    $W_{AB}(\vec{P})=mg(z_{A}-z_{B})=+mgh$ car la bille descend.

                    La figure montre que~: $h=HB=OB-OH$~; $OB=\ell$ et
                    $OH=\ell\cos\theta_0$.
                    \begin{gather*}
                        W(\vec{P})=+mg\ell(1-\cos\theta_0) \\
                        W(\vec{P})=0,250\times9,8\times1,00\times
                        (1-\cos\ang{60})\approx\qty{1.2}{\J}.
                    \end{gather*}
          \end{itemize}
\end{enumerate}
